\section{A uniform splitting estimate}\label{sec:splitting}

The analytic input is the following assertion of
\cite[Theorem~1.2]{PrimitiveRoot2}.  A finite-order Hecke character of a
number field \(F\) means a continuous finite-image character of
\(F^\times\backslash\mathbb A_F^\times\).

\begin{hypothesis}[Uniform Hecke zero-free assertion]\label{input:hecke}
Let \(F\) be a cyclotomic number field containing \(\mu_{12}\).  For every
finite-order Hecke character \(\chi\) of \(F\), the meromorphic continuation
of \(L_F(s,\chi)\) has no zeros in
\[
                         \Re s>1-10^{-6}.
\]
\end{hypothesis}

The assertion permits the pole at \(s=1\) for the principal character.
We write \(\eta=10^{-6}\) throughout this section.

\begin{proof}[Proof of Proposition~\ref{prop:splitting}]
Fix the positive prime \(q\), put \(L=\log x\), and let \(\ell\) be a prime
in the range of the proposition.  We take \(x\) large enough in terms of
\(q\) that \(\ell>q\) and \(\ell\ge5\).  Set
\[
 E_\ell=\mathbb Q(\mu_\ell),\qquad
 K=K_{\ell,q}=E_\ell(q^{1/\ell}),\qquad
 n=[K:\mathbb Q],\qquad D=|d_K|.
\]
The field \(K\) is the splitting field of \(T^\ell-q\), so it is Galois
over \(\mathbb Q\).  The cyclotomic discriminant is
\(|d_{E_\ell}|=\ell^{\ell-2}\).  Since \(q\ne\ell\), the prime \(q\) is
unramified in \(E_\ell\), and its valuation at each prime of \(E_\ell\)
above \(q\) is one.  The polynomial \(T^\ell-q\) is Eisenstein at such a
prime.  Consequently
\begin{equation}
 [K:E_\ell]=\ell,\qquad n=\ell(\ell-1).
 \label{eq:kummer-degree}
\end{equation}
The discriminant of the power basis generated by \(q^{1/\ell}\) is, up to
sign, \(\ell^\ell q^{\ell-1}\).  The relative field discriminant therefore
divides \((\ell^\ell q^{\ell-1})\mathcal O_{E_\ell}\).  The discriminant
tower formula gives
\begin{equation}
 \begin{aligned}
 D&\le \ell^{\ell(\ell-2)}
       \bigl(\ell^\ell q^{\ell-1}\bigr)^{\ell-1}
    =\ell^{\ell(2\ell-3)}q^{(\ell-1)^2},\\
 \log D+n&\ll_q n\log(2\ell).
 \end{aligned}
 \label{eq:kummer-discriminant}
\end{equation}
The same divisibility and the cyclotomic discriminant show that rational
primes outside \(q\ell\) are unramified in \(K\).

We identify the two congruences in Proposition~\ref{prop:splitting} with
complete splitting.  Let \(p\nmid q\ell\) be a prime, choose a prime
\(\mathfrak P\) of \(K\) above \(p\), and write
\(\sigma=\operatorname{Frob}_{\mathfrak P}\) for arithmetic Frobenius.
Its restriction to \(E_\ell\) sends \(\zeta_\ell\) to \(\zeta_\ell^p\),
so it is trivial exactly when \(p\equiv1\pmod\ell\).  Under this condition,
put \(\alpha=q^{1/\ell}\).  Since \(p\nmid q\), \(\alpha\) is a unit at
\(\mathfrak P\), and the defining congruence for arithmetic Frobenius gives
\[
 \frac{\sigma(\alpha)}{\alpha}\in\mu_\ell,\qquad
 \frac{\sigma(\alpha)}{\alpha}
   \equiv\alpha^{p-1}=q^{(p-1)/\ell}\pmod{\mathfrak P}.
\]
Reduction on \(\mu_\ell\) is injective at \(p\ne\ell\).  Hence
\(q^{(p-1)/\ell}\equiv1\pmod p\) is equivalent to \(\sigma\) fixing
\(\alpha\).  Together the two congruences say that \(\sigma\) fixes both
\(E_\ell\) and \(\alpha\), and hence is the identity on \(K\).  Because
\(K/\mathbb Q\) is Galois, this is precisely complete splitting of \(p\).

We transfer Input~\ref{input:hecke} to \(\zeta_K\).  Put
\[
 F_\ell=\mathbb Q(\mu_{12\ell}),\qquad
 \widetilde K=KF_\ell=F_\ell(q^{1/\ell}).
\]
Because \(F_\ell\) contains \(\mu_\ell\), the extension
\(\widetilde K/F_\ell\) is Galois and the map
\(\sigma\mapsto\sigma(q^{1/\ell})/q^{1/\ell}\) embeds its Galois group
in \(\mu_\ell\).  It is thus cyclic.  Also \(F_\ell/\mathbb Q\) is abelian
Galois, so \(\widetilde K/K\) is Galois and restriction embeds its Galois
group in \(\Gal(F_\ell/\mathbb Q)\).  It too is abelian.  Abelian Artin
formalism \cite[Chapter~VII, Corollary~(10.5) and Theorem~(10.6)]{Neukirch1999},
with the one-dimensional Artin functions identified with their associated
primitive Hecke functions, gives
\begin{equation}
 \begin{aligned}
 \zeta_{\widetilde K}(s)
  &=\prod_{\chi\in\widehat{\Gal(\widetilde K/F_\ell)}}L_{F_\ell}(s,\chi),\\
 \frac{\zeta_{\widetilde K}(s)}{\zeta_K(s)}
  &=\prod_{\substack{\chi\in\widehat{\Gal(\widetilde K/K)}\\\chi\ne1}}
       L_K(s,\chi).
 \end{aligned}
 \label{eq:abelian-zeta-factorizations}
\end{equation}
These are identities of meromorphic functions, including the local factors
at ramified primes.  The first product is zero-free on \(\Re s>1-\eta\)
by Input~\ref{input:hecke}, since \(F_\ell\) is cyclotomic and contains
\(\mu_{12}\); its principal factor has the permitted pole at one.  Each
factor of the second product is an entire nonprincipal finite-order Hecke
\(L\)-function \cite{Tate1967}.  A zero of \(\zeta_K\) in this half-plane away from one
would therefore be a zero of \(\zeta_{\widetilde K}\).  Hence
\begin{equation}
 \zeta_K(s)\ne0\qquad(\Re s>1-\eta,\ s\ne1).
 \label{eq:kummer-zero-free}
\end{equation}
The common global half-plane in the input makes this width uniform in
\(\ell\), without requiring uniform onsets for the auxiliary estimates in
its proof.  This is the abelian tower argument of
\cite[Section~9, Equation~(9.2)]{PrimitiveRoot2}, specialized to the
base \(q\).

To count completely split primes uniformly as the field varies, we use
the prime-ideal calculation from
\cite[Section~9, Equations~(9.5)--(9.9)]{PrimitiveRoot2}.  Fix a nonnegative
\(\phi\in C_c^\infty((0,\infty))\) with \(\phi\ge1\) on \([1,2]\), and put
\[
 \Phi(s)=\int_0^\infty\phi(t)t^{s-1}\dd t,\qquad
 \Psi_{K'}(x)=\sum_{\mathfrak p}\sum_{m\ge1}
       \log\Norm\mathfrak p\,
       \phi\left(\frac{(\Norm\mathfrak p)^m}{x}\right)
\]
for a number field \(K'\).  Suppose that \(K'\) has degree
\(n'=r_1+2r_2\), absolute discriminant \(D'\), and no zeta zero in
\(\Re s>1-\eta\).  Its nontrivial zeros \(\rho\) lie in
\(0<\Re\rho<1\), and in fact \(\Re\rho\le1-\eta\).  The Mellin transform
\(\Phi\) is entire and has arbitrary inverse-power decay in height on
fixed vertical strips.  Mellin inversion of the logarithmic derivative,
followed by a shift from \(\Re s=2\) to \(\Re s=-1/2\), gives
\begin{equation}
 \Psi_{K'}(x)=x\Phi(1)-\sum_\rho x^\rho\Phi(\rho)
  -(r_1+r_2-1)\Phi(0)
  +O_\phi\left(x^{-1/2}(\log D'+n')\right).
 \label{eq:smooth-prime-ideal-formula}
\end{equation}
Here zeros are counted with multiplicity.  The term at the origin comes
from the zero of \(\zeta_{K'}\) there of order \(r_1+r_2-1\); it is absent
when that order is zero.  There are no other trivial zeros between the two
lines.  To see the uniformity of the error, the functional equation gives
on \(\Re s=-1/2\)
\[
 \left|\frac{\zeta_{K'}'}{\zeta_{K'}}(s)\right|
 \ll \log D'+n'\log(2+|\Im s|).
\]
Indeed the Euler series at \(1-s\), with real part \(3/2\), is \(O(n')\),
and the logarithmic derivatives of the archimedean factors contribute
\(O(n'\log(2+|\Im s|))\).  The rapid decay of \(\Phi\) gives the error in
Equation~\eqref{eq:smooth-prime-ideal-formula}.  The usual rectangular
contours may be taken with horizontal edges away from zeros.  The zero count
below supplies such edges, and the Hadamard partial-fraction formula bounds
the logarithmic derivative polynomially on them for each fixed \(K'\).
The decay of \(\Phi\) then makes their integrals vanish; the same zero count
makes the zero sum absolutely convergent.

For the required zero count, let \(N_{K'}(H)\) count nontrivial zeros with
\(|\Im\rho|\le H\), including multiplicity.  The uniform formula of
\cite[Corollary~1.2, Equation~(1.3)]{HasanalizadeShenWong2022} is
\[
 N_{K'}(H)=\frac H\pi
   \log\left(D'\left(\frac{H}{2\pi e}\right)^{n'}\right)
   +O\bigl(\log D'+n'\log H+n'\bigr)\qquad(H\ge1).
\]
For \(|T|\ge2\), the unit band is bounded by
\(N_{K'}(|T|+2)-N_{K'}(|T|-1)\).  The main term changes by
\(O(\log D'+n'\log(2+|T|))\) over this interval, and both errors have the
same bound.  For \(|T|<2\), use \(N_{K'}(3)\).  Thus
\begin{equation}
 \#\{\rho:T\le\Im\rho<T+1\}
 \ll \log D'+n'\log(2+|T|),
 \label{eq:unit-height-zero-count}
\end{equation}
uniformly in \(K'\) and \(T\).  Using it with the decay of \(\Phi\) on
\(0\le\Re s\le1\) gives, for \(x\ge1\),
\begin{equation}
 \begin{aligned}
 \sum_\rho|x^\rho\Phi(\rho)|
 &\ll_\phi x^{1-\eta}\sum_{j\in\mathbb Z}(1+|j|)^{-3}
       \bigl(\log D'+n'\log(2+|j|)\bigr)\\
 &\ll_\phi x^{1-\eta}(\log D'+n').
 \end{aligned}
 \label{eq:smooth-zero-sum}
\end{equation}
The origin term in Equation~\eqref{eq:smooth-prime-ideal-formula} is
nonpositive, since \(\Phi(0)\ge0\).  Equations
\eqref{eq:smooth-prime-ideal-formula} and \eqref{eq:smooth-zero-sum} thus
imply the uniform upper bound
\begin{equation}
 \Psi_{K'}(x)\ll_\phi x+x^{1-\eta}(\log D'+n')\qquad(x\ge1).
 \label{eq:smooth-prime-ideal-upper}
\end{equation}
In particular, the implied constant depends only on the fixed test function,
and not on the field.

Apply Equation~\eqref{eq:smooth-prime-ideal-upper} to \(K'=K\), using
Equation~\eqref{eq:kummer-zero-free}.  Every rational prime \(p\) splitting
completely in \(K\) supplies \(n\) prime ideals of norm \(p\).  When
\(x<p<2x\), their first-power terms contribute at least \(n\log p>nL\)
to \(\Psi_K(x)\).  Equations~\eqref{eq:kummer-degree} and
\eqref{eq:kummer-discriminant} therefore give
\begin{equation}
 \begin{aligned}
 &\#\{p\text{ prime}:x<p<2x,\ p\text{ splits completely in }K\}\\
 &\qquad\ll_q \frac{x}{\ell(\ell-1)L}
       +\frac{x^{1-\eta}\log(2\ell)}{L}
 \ll_q \frac{x}{\ell(\ell-1)L}+x^{1-\eta}.
 \end{aligned}
 \label{eq:complete-splitting-bound}
\end{equation}
The last inequality uses \(\log\ell<L^{0.3}\) and sufficiently large \(x\).
By the Frobenius identification above,
Equation~\eqref{eq:complete-splitting-bound} proves the proposition.
\end{proof}
